\clearpage
\subsection{测试结果与缺陷分析}
\label{sec:impl-19-results}
\par\medskip\noindent\textbf{本轮结果汇总。}
2026年11月9日，Windows与Node.js v24.13.0下，16个\DesignCode{*.test.mjs}文件共执行141项测试，141项通过，失败、取消、跳过及待办均为零。框架记录总耗时2606.1751毫秒，约2.606秒；这是测试套件运行时间，不是单次HTTP响应时间或模型推理时间。

三个浏览器脚本退出码均为0。证据采集器测得的墙钟时间包括进程启动，与框架内部计时不同；不能把两者取较小值作为性能结论。历史120项结果保留在原目录，新增Qwen等测试后项数不同，不代表产品需求数量增加了相同数目。

\begin{table}[H]\centering\small
\caption{缺陷、修复与验证链示例}
\begin{tabularx}{\linewidth}{@{}p{33mm}X@{}}\toprule
问题 & 修复思路与当前证据 \\ \midrule
合法标签夹带未知兴趣即失败 & 过滤后保留合法值；本轮解析、重试及数据库回归通过 \\
HTTP环境生图编号函数缺失 & 使用兼容编号生成；本轮浏览器模拟缺失函数仍能生成与保存 \\
旧情绪或风格覆盖新消息 & 返回时核对消息与风格；本轮延迟返回竞争用例通过 \\
头像与聊天立绘混用 & 独立存储和接口；媒体分离与迁移回归通过 \\ \bottomrule
\end{tabularx}\end{table}

\par\medskip\noindent\textbf{不能从通过结果推出的结论。}
全部断言通过并不排除未覆盖分支、长期并发、真实模型变更或新的输入组合。本轮没有重测历史16组性能数据，也没有给出公网承载量、长期可用性和情绪准确率。模型参考分的在途关闭等风险应单独安排异步授权撤回测试，不能被普通开关断言抵消。

本章保留失败分支与待验证边界，是为了让后续补测能够定位到具体假设，而不是将“141项通过”扩大成“所有功能完全验收”。原始记录随源码提供，可复查执行范围和退出状态。
